\usepackage{amssymb}
\usepackage{colortbl}
\usepackage{eso-pic}
\usepackage{fancyhdr}
\usepackage{lastpage}
\usepackage{needspace}
\usepackage{pdflscape}

\definecolor{RTTableHeader}{gray}{0.90}
\definecolor{RTTableStripe}{gray}{0.96}
\arrayrulecolor{black!30}
\renewcommand{\arraystretch}{1.10}

\providecommand{\rtclassification}{UNCLASSIFIED}
\newcommand{\rtpagecount}{\thepage/\pageref*{LastPage}}
\newcommand{\rtconfigurepagestyle}{%
  \fancyhf{}%
  \fancyhead[C]{\textbf{\footnotesize\rtclassification}}%
  \fancyfoot[C]{%
    \begin{tabular}{@{}c@{}}
      \textbf{\footnotesize\rtclassification}\\[-1pt]
      \scriptsize\rtpagecount
    \end{tabular}}%
  \renewcommand{\headrulewidth}{0.4pt}%
  \renewcommand{\footrulewidth}{0.4pt}%
}
\fancypagestyle{rtclassified}{\rtconfigurepagestyle}
\fancypagestyle{plain}{\rtconfigurepagestyle}
\fancypagestyle{empty}{\rtconfigurepagestyle}
\pagestyle{rtclassified}
\setlength{\headheight}{16pt}
\setlength{\footskip}{32pt}

% pdflscape rotates the page after fancyhdr has placed its marks. Wide-table
% pages therefore receive marks at ship-out in the coordinates that become the
% visual top and bottom after rotation.
\fancypagestyle{rtlandscapepage}{%
  \fancyhf{}%
  \renewcommand{\headrulewidth}{0pt}%
  \renewcommand{\footrulewidth}{0pt}%
}
\newcommand{\rtlandscapemarks}{%
  \AddToShipoutPictureFG{%
    \AtPageLowerLeft{%
      \put(\LenToUnit{\dimexpr\paperwidth-20pt\relax},\LenToUnit{0.5\paperheight}){%
        \makebox(0,0){\rotatebox{90}{\textbf{\footnotesize\rtclassification}}}}%
      \put(\LenToUnit{\dimexpr\paperwidth-32pt\relax},\LenToUnit{2.5cm}){%
        \rule{0.4pt}{\dimexpr\paperheight-5cm\relax}}%
      \put(\LenToUnit{32pt},\LenToUnit{2.5cm}){%
        \rule{0.4pt}{\dimexpr\paperheight-5cm\relax}}%
      \put(\LenToUnit{20pt},\LenToUnit{0.5\paperheight}){%
        \makebox(0,0){\rotatebox{90}{\textbf{\footnotesize\rtclassification}}}}%
      \put(\LenToUnit{\dimexpr\paperwidth-8pt\relax},\LenToUnit{0.5\paperheight}){%
        \makebox(0,0){\rotatebox{90}{\scriptsize\rtpagecount}}}%
    }%
  }%
}
\newenvironment{rtlandscape}{%
  \begin{landscape}%
  \pagestyle{rtlandscapepage}%
  \thispagestyle{rtlandscapepage}%
  \rtlandscapemarks%
}{%
  \end{landscape}%
  \ClearShipoutPictureFG%
  \pagestyle{rtclassified}%
}

\newcommand{\rtstarttemplate}{%
  \ifdim\pagetotal>0pt\relax
    \clearpage
  \fi
}

% Pandoc calls \subtitle after this header is loaded. Capture it separately so
% the cover can give the title and subtitle distinct typographic treatment.
\newcommand{\rtsubtitle}{}
\providecommand{\subtitle}[1]{\renewcommand{\rtsubtitle}{#1}}

\makeatletter
\renewcommand{\maketitle}{%
  \thispagestyle{rtclassified}%
  \vspace*{0.18\textheight}%
  \begin{center}
    {\fontsize{25}{30}\selectfont\bfseries \@title\par}
    \vspace{1.2cm}
    \rule{0.70\textwidth}{0.8pt}\par
    \vspace{1.2cm}
    {\large \rtsubtitle\par}
    \vspace{1.4cm}
    {\normalsize\bfseries \@author\par}
    \vspace{0.45cm}
    {\large\bfseries \@date\par}
  \end{center}
}
\makeatother

% The body starts after an intentionally blank classified leaf following the
% generated contents pages.
\let\rtoriginaltableofcontents\tableofcontents
\renewcommand{\tableofcontents}{%
  \rtoriginaltableofcontents
  \clearpage
  \null
  \thispagestyle{rtclassified}
  \clearpage
}

% Unicode mappings required by the Markdown source when using pdflatex.
\DeclareUnicodeCharacter{00B7}{\ensuremath{\cdot}}
\DeclareUnicodeCharacter{00F7}{\ensuremath{\div}}
\DeclareUnicodeCharacter{2026}{\ldots}
\DeclareUnicodeCharacter{2191}{\ensuremath{\uparrow}}
\DeclareUnicodeCharacter{2193}{\ensuremath{\downarrow}}
\DeclareUnicodeCharacter{2194}{\ensuremath{\leftrightarrow}}
\DeclareUnicodeCharacter{2610}{\ensuremath{\square}}
\DeclareUnicodeCharacter{25A3}{\ensuremath{\boxdot}}
\DeclareUnicodeCharacter{2714}{\ensuremath{\checkmark}}
